\documentclass[]{article}

\usepackage{amsmath}
\usepackage{multirow}
\usepackage{natbib}
\usepackage{graphicx} 


\begin{document}

% Auto-generated figures file
% \usepackage{graphicx}

\begin{figure}[h!]
    \centering
    \includegraphics[width=0.5\textwidth]{../input_files/plots/step_1_fit_summary_1_1775343995.png}
    \caption{Representative results of the Maximum Likelihood Estimation (MLE) framework applied to four underdamped harmonic oscillators. Each panel displays the noisy time-series displacement data (grey points) alongside the best-fit model trajectory (red line) obtained via non-linear least-squares minimization. The plots visually confirm the high fidelity of the parameter recovery, showing that the model accurately captures the oscillatory dynamics and exponential decay over the full 20-second interval for systems with varying damping ratios and signal-to-noise levels.}
    \label{fig:image0}
\end{figure}

\begin{figure}[h!]
    \centering
    \includegraphics[width=0.5\textwidth]{../input_files/plots/step_2_diagnostic_plots_1775344014.png}
    \caption{Performance of the Maximum Likelihood Estimation framework for recovering oscillator parameters. Parity plots (top row) for the damping rate ($\gamma$) and natural frequency ($\omega$) show excellent agreement between fitted and true values, with points clustering tightly along the identity line. The bottom row shows the relative error in the estimation of $\gamma$. The error exhibits an inverse relationship with the Signal-to-Noise Ratio (SNR) (bottom left) and a slight increase as the damping ratio increases (bottom right), indicating a degradation in precision for more rapidly decaying signals.}
    \label{fig:image1}
\end{figure}



\end{document}
                